\section{Methods}\label{sec:methods}

\subsection{PIGEAN Model}
\hypertarget{pigeanpior}{The} PIGEAN model specifies that (a) a subset of genes are ``relevant" to a trait and (b) the probability a gene is relevant to a trait is affected (additively) by the gene's annotations. Specifically, for an arbitrary trait $T$ and genes $g_i\ldots g_N$, let $d_i \in\{0,1\}$ indicate whether gene $i$ is involved in $T$ and let $y_i=\Pr(d_i=1)$. Given a set of (binary or real valued) annotations $X_{ij}$ (where $j=1\dots M$) for $g_i$, we then model:
\[
\log\frac{y_i}{1-y_i} = \sum_j X_{ij}\beta_j + \alpha_i + \log\frac{\pi}{1-\pi} 
\]
\begin{align*}
\beta_j \sim \delta_j\mathbf{N}(0, \sigma_j^2) \quad
\alpha_i = \alpha\sum_j X_{ij} \quad
\sigma_j^2 = \sigma^2\text{Var}_i(X_{ij})^w \quad
\delta_j \sim \text{Bern}(p) \quad
\end{align*}
where $\pi$ is the prior probability that a gene is involved in disease, which we typically set to 0.05 ($\sim 1000$ genes) for complex traits and 0.01 ($\sim 200$ genes) for rare diseases, and the four parameters ($\alpha$, $\sigma$, $w$, and $p$) enforce regularization: $\sigma$ and $w$ penalize high variance (i.e. "large") genes and annotations, $\sigma$ shrinks gene annotation effect sizes, and $p$ enforces sparsity of gene annotation effect sizes. As described in the Online Methods, we learn $p$ from the input data for trait $T$, set $w$ and $\alpha$ to limit dependence of estimands on the variance of annotations or genes, and set $\sigma$ to an empirically determined large trait-independent value (learning trait-specific $\sigma$ values is possible but, empirically, has little impact on results and introduces the potential for less robustness during model application across many traits).

We do not know $d_i$ or $\beta_j$ (they are latent variables), but we can jointly infer them conditional upon observed genetic association statistics $\mathcal{S}$. By employing Gibbs sampling, we can sample from this distribution $\prob(d, \beta \mid \mathcal{S})$ if we can sample from $\prob(d \mid \beta, \mathcal{S})$ and $\prob(\beta \mid d, \mathcal{S})=\prob(\beta \mid d)$. We can apply previous ideas[] to sample from $\prob(\beta \mid d)$ (Methods), and we consider four types of genetic association statistics of increasing complexity for calculating the conditional probability $\prob(d \mid \beta, \mathcal{S})$: rare diseases in which a subset of disease-relevant genes are "known", WES gene-level burden test statistics, genome-wide SNP summary statistics from a GWAS, and a list of only significant SNPs from a GWAS. More complicated scenarios including multiple types of summary statistics (e.g. GWAS summary statistics alongside WES burden tests of ultra-rare variants) can be handled in a straightforward manner, provided the summary statistics are independent (see Methods).

In the rare disease scenario, we assume that we are given $\prob(d \mid \mathcal{S})$, either fixed to a high value for a small number of "known" genes or estimated from qualitative confidence "tiers"[].  It is straightforward (Methods) to show that the log of the conditional odds of $d_i$ is the sum of its prior odds conditional upon $\beta$ and the unconditional gene bayes factor (which can be calculated from $\prob(d \mid \mathcal{S})$ and $\pi$), under the assumption that gene annotations included in the model are not influenced by knowledge of disease genes for $T$ (we investigate the feasibility of this assumption and the consequences of its violation below). In the WES scenario, the calculations are identical except, rather than being given $\prob(d \mid \mathcal{S})$ and converting it to gene bayes factors, we estimate the gene bayes factors from the burden effect sizes and standard errors using an approach previously developed for SNPs (Methods).

In the genome-wide GWAS summary statistics scenario, we calculate $\prob(d \mid \beta, \mathcal{S})$ using an extension of the HuGE calculator, a previously proposed set of simple rules to estimate gene-level probabilities of disease relevance conditional upon nearby SNP associations (Methods). The extensions we use within PIGEAN (referred to as HuGEX; see Methods) address some of the obvious limitations of the HuGE calculator while preserving its advantages (simplicity, speed, and interpretability). In the  scenario of partial GWAS summary statistics, we apply the HuGEX model to all genes nearby SNPs with summary statistics available, and set the remaining genes to have $\prob(d \mid \beta, \mathcal{S})=\prob(d \mid \beta)$. 

The joint posterior distribution $\Pr(d,\beta|\mathcal{S})$ inferred by PIGEAN can be distilled into four sets of summary statistics (Methods): $E[\beta_j]$, the effect size for the $j^{th}$  gene annotation; $ABF^I_i$, the "indirect" approximate Bayes Factor (ABF) that represents evidence from the annotations for gene $g_i$; $ABF^D_i$, the "direct" ABF that represents evidence from genetic associations for gene $g_i$; and $ABF_i$, the combined $ABF$ for gene $g_i$. The overall probability of disease relevance for a gene can then be expressed straightforwardly as a function of $ABF_i$ and the prior $\pi$. 

These equations offer three interpretations of the PIGEAN model (Methods). In the "principled priors" interpretation, the posterior probability of a gene's disease relevance is estimated using genetic associations as data ($ABF_i^D$) and the annotations to adjust the prior (a combination of $ABF_i^I$ and $\pi$), offering bayesian approach to use both genetic associations and gene annotations to prioritize genes. In the "bottom-up" interpretation, PIGEAN generalizes the popular PoPS approach; it can be shown (Methods) that if MAGMA Z-scores are used to estimate the log-odds of $ABF_i^D$, sparsity constraints on the $\beta_j$ values are replaced by L2 regularization, and full Gibbs sampling from the posterior is replaced by a single regression, then PIGEAN reduces to PoPS --- these alterations elucidate the quantity PoPS  estimates as well as the approximations in inference and data that it makes. Finally, the "indirect genetic support" interpretation of PIGEAN demonstrates that genetic support is due not only to genetic associations nearby a gene ($ABF_i^D$) but also to genetic associations nearby other genes ($ABF_i^I$) that indirectly support the gene through shared annotations.

We implemented PIGEAN as a Python program and applied it to {{NUM_GWAS_FULL_SUMSTATS}} full summary statistic GWAS datasets from the Association to Function Knowledge Portal, lists of reported significant GWAS associations for {{NUM_TRAITS_CURATED}} traits from the GWAS catalog, and {{NUM_RARE_DISEASE_GROUPS:,}} rare disease gene lists from Orphanet (\textbf{Supplementary Table 1}). For validation purposes, we also analyzed 20 simulated "null" GWAS in which all SNPs had zero true effect sizes, and 10 simulated GWAS in which SNPs had true effect sizes consistent with the PIGEAN model (Methods). We considered a total of {{NUM_TOTAL_GENE_ANNOTATIONS:,}} gene annotations, including (Supplementary Table 2) knockout mouse phenotypes ({{NUM_MPO_TERMS:,}} Mammalian Phenotype Ontology terms), membership in gene sets and pathways ({{NUM_MSIGDB_GENESETS:,}} gene sets from the Molecular Signatures Database), occurrence in a PubMedCentral manuscript figure ({{NUM_PFOCR_SETS:,}} gene sets from the Pathway Figure OCR project), co-occurrence with terms in PubMed abstracts ({{NUM_MESH_TERMS:,}} MeSH terms), network interaction partners ({{NUM_STRING_NEIGHBORS:,}} gene neighbors from STRING), and single-cell gene expression patterns ({{NUM_SINGLE_CELL_CLUSTERS:,}} experiment clusters curated as part of the PoPS method).

\subsection{Validation on Simulations}
As a negative control, we first applied PIGEAN to the 20 "null" simulated GWAS and found that, as expected, inferred gene annotation effect sizes were all close to zero (Extended Data Figure~\ref{fig:null_sim}a) as were indirect, direct, and combined gene ABFs (Supplementary Figure \ref{fig:null_sim}bcd). Similarly, when we applied PIGEAN to the 50 real GWAS but randomly permuted genes within the 166,868 gene annotations, we also observed annotation effect sizes close to zero (Supplementary  Figure \ref{fig:null_sim}e). Across our 20 null simulations and 166,868 annotations only one gene had combined gene ABF above 3 (Supplementary Figure \ref{fig:null_sim}g), with the dominant contribution coming from direct ABFs (Supplementary Figure \ref{fig:null_sim}h), and across the 50 real GWAS with permuted annotations, only one annotation had effect size above 0.004 (Supplementary Figure \ref{fig:null_sim}f). Based on these results, annotation effect sizes of 0.005 and gene combined ABFs of 3 roughly correspond to "study-wide significant" p-value thresholds (Supplementary Table 3).

We also used simulations with non-zero annotation effect sizes to understand the behavior of PIGEAN in response to different model parameters. Accuracy of PIGEAN is affected by noise from two sources: the incomplete ability of annotations to predict disease-relevant genes ("annotation noise"), and the incomplete ability of summary statistics to identify disease-relevant genes ("association noise"). Focusing first on annotation noise, we ran PIGEAN with a truncated model in which we treated each gene's probability of disease relevance ($d_i$) as observed rather than inferred from summary statistics. These simulations yielded four main findings. First, PIGEAN in general conservatively estimates both annotation effect sizes (Supplementary Figure \ref{fig:non_null_sim}a) and gene ABFs (Supplementary Figure \ref{fig:non_null_sim}b), due to its "spike and slab" regularization that sets estimated effect sizes to zero and shrinks estimated effect sizes toward zero. Second, annotations with negative effects (i.e. that reduce genetic support) were shrunk substantially more than those with positive effects, as (consequently) were genes with ABFs < 1, an asymmetry that naturally results from the higher standard errors for negative gene set effect sizes when there are only a small number of disease-relevant genes (Supplementary Figure \ref{fig:non_null_sim}cd). Third, PIGEAN's accuracy was also affected by simulation parameters, with more accurate inferences possible for larger simulated effect size magnitudes (larger $\sigma$ Supplementary Figure \ref{fig:sim_params}a) and (to a lesser extent) more non-zero simulated effect sizes (larger $p$; Supplementary Figure \ref{fig:sim_params}b). Additional noise added to the simulated gene ABFs affected accuracy of the inferred ABFs but (perhaps non-intuitively) did not affect accuracy of the inferred annotation effect sizes (Supplementary Figure \ref{fig:sim_params}c). Finally, due to the nature of Gibbs sampling, estimates produced by PIGEAN varied across runs, but with sufficient numbers of iterations and chains estimates were stable (Supplementary Figure \ref{fig:sim_stability}).

These simulations confirmed the importance PIGEAN's joint inference and regularization features. First, its full joint inference procedure (in which gene and gene set values are iteratively re-estimated) increased accuracy compared to a model in which gene set effect sizes were only estimated once (Supplementary Figure \ref{fig:sim_onestep}). Second, even though simulated gene ABFs had no dependence on the number of annotations for the gene, regularization (via its  parameter) to penalize genes with many annotations was necessary to avoid biases and ABF accuracy reductions induced (as described in Supplementary Figure \ref{fig:non_null_sim}cd) by asymmetric annotation standard errors (Supplementary Figure \ref{fig:sim_induced}). Similarly, even though simulated annotation effect sizes had no dependence on the number of genes with the annotation, regularization (via its w parameter) to penalize "large" annotations was necessary to avoid biases induced by the smaller standard errors of larger annotations (Supplementary Figure \ref{fig:sim_reg}; Methods).

For simulations with a small number of non-zero annotation effects, PIGEAN was able to infer reasonably accurate estimates of its sparsity parameter ($p$) (Supplementary Figure \ref{fig:sim_param_infer}), although it significantly underestimated $p$ as the fraction of non-zero effects increased. Our attempts to estimate PIGEAN's shrinkage parameter ($\sigma$) for individual simulations suggested that there are typically too few annotations  to estimate this parameter accurately, and we therefore tuned this parameter empirically (Methods). Fortunately, varying these parameters from their simulated values by even an order of magnitude did not dramatically reduce PIGEAN's estimation accuracy (Supplementary Figures \ref{fig:sim_vary_sigma} and \ref{fig:sim_vary_p}).

In addition to these simulations focused exclusively on "annotation noise", we also evaluated the impact of  "association noise" on PIGEAN's performance -- that is, the reality that (due to finite sample sizes and the vagaries of evolution) many disease-relevant genes will not lie near detectable association signals. We studied association noise in the context of GWAS summary statistics, given the well-documented difficulties in inferring effector genes from these data. For each of three annotation noise settings (Methods), we simulated five GWAS of increasing sample sizes (using linkage disequilibrium from the 1000 Genomes Project, SNP to gene linkages from a previously published method, and SNP effect size parameters from previously published simulations; Methods). As expected, the use of GWAS data rather than observed disease probabilities degraded the accuracy of PIGEAN (Supplementary Figure \ref{fig:sim_gwas_gene_set}a). With perfect knowledge of SNP to gene linkages Supplementary Figure \ref{fig:sim_gwas_gene_set}b), less annotation noise Supplementary Figure \ref{fig:sim_gwas_gene_set}c), or larger sample sizes Supplementary Figure \ref{fig:sim_gwas_gene_set}d), the accuracy of PIGEAN increased, but only to limited extent. Accuracy was much higher for marginal annotation effect sizes than for joint effect sizes Supplementary Figure \ref{fig:sim_gwas_gene_set}e), which is a common feature of multivariate models with correlated predictors. When PIGEAN inferred its regularization parameters, rather than being given the true parameters used in simulations, accuracy degraded further Supplementary Figure \ref{fig:sim_gwas_gene_set}f). Notably, however, in all scenarios PIGEAN conservatively estimated annotation effect sizes, suggesting that the overwhelming impact of GWAS annotation noise is to decrease power rather than to increase false positives.

We conclude that, under idealized simulations, PIGEAN behaves broadly as expected. Three noteworthy properties are (a) the difficulty of inferring negative genetic support (due to known properties of GWAS but perhaps unappreciated properties of models with few disease-relevant genes); (b) the importance of regularization and iterative inference in gene set analysis, and (c) its conservative nature when confronted with annotation noise and association noise.

\subsection{Validation on Real Data}
\paragraph{Selection of Validation Phenotypes}
%TODO: Also include bottom line and/or credible set information for each phenotype to prove we looked at phenos with varying GWAS significance levels
For common traits, we selected 13 phenotypes (Table~\ref{tab:common-validation-traits}) encompassing both continuous and dichotomous traits with GWAS data of varying significance but which we had curated gold-standard genesets for validation. This diversity allows us to evaluate PIGEAN and respective PIGEAN models under a broad range of real-world conditions. For rare traits, we chose 40 phenotypes (Table~\ref{tab:rare-validation-traits}) that span multiple organ systems, inheritance patterns, and pathophysiological mechanisms. This set includes both specific rare diseases (with more than 10 associated genes) and groups of disorders representing higher-level disease clusters. To ensure a balanced representation without excessive overlap in gene lists, we verified that the maximum Jaccard similarity among groups of disorders is 0.14 and 0.11 among specific diseases, with average similarities of 0.00653 and 0.0016, respectively. By covering a wide range of disease etiologies and genetic architectures, we can thoroughly assess our method's performance across diverse traits.

\begin{table}[ht]
\centering
\footnotesize
\begin{tabular}{l c c}
\hline
\textbf{Disease Name} & \textbf{\# Sig GWAS Loci} & \textbf{\# Credible Sets} \\
\hline
\href{https://a2f.hugeamp.org/phenotype.html?phenotype=AF}{Atrial Fibrillation} & TBD & TBD \\
\href{https://a2f.hugeamp.org/phenotype.html?phenotype=ALLOA}{Osteoarthritis} & TBD & TBD \\
\href{https://a2f.hugeamp.org/phenotype.html?phenotype=AnyCVD}{Cardiovascular Disease} & TBD & TBD \\
\href{https://a2f.hugeamp.org/phenotype.html?phenotype=Eczema}{Eczema} & TBD & TBD \\
\href{https://a2f.hugeamp.org/phenotype.html?phenotype=HDL}{High-Density Lipoprotein Cholesterol (HDL)} & TBD & TBD \\
\href{https://a2f.hugeamp.org/phenotype.html?phenotype=HF}{Heart Failure} & TBD & TBD \\
\href{https://a2f.hugeamp.org/phenotype.html?phenotype=HYPERTENSION}{Hypertension} & TBD & TBD \\
\href{https://a2f.hugeamp.org/phenotype.html?phenotype=LDL}{Low-Density Lipoprotein Cholesterol (LDL)} & TBD & TBD \\
\href{https://a2f.hugeamp.org/phenotype.html?phenotype=MDD}{Major Depressive Disorder} & TBD & TBD \\
\href{https://a2f.hugeamp.org/phenotype.html?phenotype=MI}{Myocardial Infarction} & TBD & TBD \\
\href{https://a2f.hugeamp.org/phenotype.html?phenotype=RhA}{Rheumatoid Arthritis} & TBD & TBD \\
\href{https://a2f.hugeamp.org/phenotype.html?phenotype=SCZ}{Schizophrenia} & TBD & TBD \\
\href{https://a2f.hugeamp.org/phenotype.html?phenotype=T2D}{Type 2 Diabetes} & TBD & TBD \\
\hline
\end{tabular}
\caption{Common Disease Validations with hyperlinked names and placeholders for GWAS loci and credible sets.}
\label{tab:common-validation-traits}
\end{table}


\begin{table}[ht]
\centering
\footnotesize
\begin{tabular}{l l l l c}
\hline
\textbf{Disease Name} & \textbf{Orphanet ID} & \textbf{Type} & \textbf{\# Genes} \\
\hline
Neuromuscular disease & \href{https://www.orpha.net/en/disease/detail/68381}{68381} & Group & 374 \\
Inherited arrhythmogenic cardiomyopathy & \href{https://www.orpha.net/en/disease/detail/247}{247} & Group & 13 \\
Primary immunodeficiency & \href{https://www.orpha.net/en/disease/detail/101997}{101997} & Group & 328 \\
Mucopolysaccharidosis & \href{https://www.orpha.net/en/disease/detail/79213}{79213} & Group & 11 \\
Sphingolipidosis & \href{https://www.orpha.net/en/disease/detail/79225}{79225} & Group & 17 \\
Ciliopathy & \href{https://www.orpha.net/en/disease/detail/363250}{363250} & Group & 282 \\
Rare sleep disorder & \href{https://www.orpha.net/en/disease/detail/68354}{68354} & Group & 14 \\
Rare pervasive developmental disorder & \href{https://www.orpha.net/en/disease/detail/168778}{168778} & Group & 25 \\
Sexchromosome anomaly & \href{https://www.orpha.net/en/disease/detail/98155}{98155} & Group & 21 \\
Autosomal trisomy & \href{https://www.orpha.net/en/disease/detail/98130}{98130} & Group & 16 \\
Hereditary ataxia & \href{https://www.orpha.net/en/disease/detail/183518}{183518} & Group & 182 \\
Peroxisomal disease & \href{https://www.orpha.net/en/disease/detail/68373}{68373} & Group & 28 \\
Mitochondrial disease & \href{https://www.orpha.net/en/disease/detail/68380}{68380} & Group & 208 \\
Disorder of purine or pyrimidine metabolism & \href{https://www.orpha.net/en/disease/detail/79224}{79224} & Group & 50 \\
Congenital myasthenic syndrome & \href{https://www.orpha.net/en/disease/detail/590}{590} & Group & 27 \\
Rare hypertrophic cardiomyopathy & \href{https://www.orpha.net/en/disease/detail/217569}{217569} & Group & 81 \\
Amyloidosis &\href{https://www.orpha.net/en/disease/detail/69}{69} & Group & 11 \\
Hemoglobinopathy & \href{https://www.orpha.net/en/disease/detail/68364}{68364} & Group & 12 \\
Spondyloepiphyseal dysplasia & \href{https://www.orpha.net/en/disease/detail/253}{253} & Group & 38 \\
Ectodermal dysplasia syndrome & \href{https://www.orpha.net/en/disease/detail/79373}{79373} & Group & 99 \\
Amyotrophic lateral sclerosis & \href{https://www.orpha.net/en/disease/detail/803}{803} & Specific & 37 \\
Zellweger syndrome & \href{https://www.orpha.net/en/disease/detail/912}{912} & Specific & 13 \\
Fanconi anemia & \href{https://www.orpha.net/en/disease/detail/84}{84} & Specific & 22 \\
RomanoWard syndrome & \href{https://www.orpha.net/en/disease/detail/101016}{101016} & Specific & 19 \\
Cystic fibrosis & \href{https://www.orpha.net/en/disease/detail/586}{586} & Specific & 19 \\
Noonan syndrome & \href{https://www.orpha.net/en/disease/detail/648}{648} & Specific & 14 \\
Acute promyelocytic leukemia & \href{https://www.orpha.net/en/disease/detail/520}{520} & Specific & 13 \\
Bulbospinal muscular atrophy & \href{https://www.orpha.net/en/disease/detail/206701}{206701} & Specific & 16 \\
Glioblastoma & \href{https://www.orpha.net/en/disease/detail/360}{360} & Specific & 12 \\
Williams syndrome & \href{https://www.orpha.net/en/disease/detail/904}{904} & Specific & 20 \\
Kallmann syndrome  & \href{https://www.orpha.net/en/disease/detail/478}{478} & Specific & 21 \\
Hirschsprung disease & \href{https://www.orpha.net/en/disease/detail/388}{388} & Specific & 14 \\
Brugada syndrome  & \href{https://www.orpha.net/en/disease/detail/130}{130} & Specific & 22 \\
Dyskeratosis congenita & \href{https://www.orpha.net/en/disease/detail/1775}{1775} & Specific & 13 \\
Behet disease & \href{https://www.orpha.net/en/disease/detail/117}{117} & Specific & 15 \\
Jeune syndrome & \href{https://www.orpha.net/en/disease/detail/474}{474} & Specific & 11 \\
Congenital hypothyroidism & \href{https://www.orpha.net/en/disease/detail/442}{442} & Specific & 28 \\
DiamondBlackfan anemia & \href{https://www.orpha.net/en/disease/detail/124}{124} & Specific & 26 \\
Leber hereditary optic neuropathy & \href{https://www.orpha.net/en/disease/detail/104}{104} & Specific & 12 \\
WalkerWarburg syndrome & \href{https://www.orpha.net/en/disease/detail/899}{899} & Specific & 14 \\
\hline
\end{tabular}
\caption{Validation sets containing 40 diseases (20 groups of disorders and 20 specific disorders), with Orphanet links and placeholders for gene counts.}
\label{tab:rare-validation-traits}
\end{table}

\begin{table}[htbp]
\centering
\caption{LOCO validation results across traits}
\label{tab:loco-validation}
\begin{tabular}{lrrrcc}
\toprule
Trait & Pearson $r$ & Spearman $\rho$ & Beta & P-value & Adj. P-value \\
\midrule
AF & 0.270 & 0.132 & 0.258 & $7.27 \times 10^{-31}$ & $<0.001$ \\
ALLOA & 0.196 & 0.077 & 0.184 & $1.17 \times 10^{-46}$ & $<0.001$ \\
Eczema & 0.294 & 0.150 & 0.280 & $3.28 \times 10^{-47}$ & $<0.001$ \\
HDL & 0.296 & 0.029 & 0.289 & $3.71 \times 10^{-30}$ & $<0.001$ \\
HF & 0.144 & 0.048 & 0.134 & $1.21 \times 10^{-22}$ & $<0.001$ \\
HYPERTENSION & 0.186 & 0.081 & 0.172 & $1.84 \times 10^{-26}$ & $<0.001$ \\
LDL & 0.258 & 0.016 & 0.251 & $1.91 \times 10^{-22}$ & $<0.001$ \\
MDD & 0.044 & 0.033 & 0.034 & $4.31 \times 10^{-5}$ & $<0.001$ \\
MI & 0.196 & 0.048 & 0.183 & $7.19 \times 10^{-18}$ & $<0.001$ \\
RhA & 0.346 & 0.159 & 0.331 & $5.59 \times 10^{-53}$ & $<0.001$ \\
SCZ & 0.160 & 0.065 & 0.145 & $3.83 \times 10^{-26}$ & $<0.001$ \\
T2D & 0.228 & 0.063 & 0.218 & $1.04 \times 10^{-22}$ & $<0.001$ \\
\bottomrule
\end{tabular}
\begin{tablenotes}
\small
\item Pearson $r$ = Pearson correlation coefficient; Spearman $\rho$ = Spearman rank correlation coefficient; Beta = regression coefficient from adjusted model; P-value = p-value from adjusted regression model; Adj. P-value = Benjamini-Hochberg adjusted two-sided p-value.
\end{tablenotes}
\end{table}


\subsection{Gene Set Library Curation}
% Section about MeSH Text Mining models
For text mining models, we curated text mining based gene sets from three libraries: Gene-Mesh Co-Occurance via Pubmed publications, String DB, and RummaGene. For the Gene-MeSH co-occurance gene sets, we downloaded Gene to PubMed article mappings from the PubTator3 \cite{wei_pubtator_2024} FTP server\footnote{The gene2pubtator3.gz file is available at the following \href{https://ftp.ncbi.nlm.nih.gov/pub/lu/PubTator3/}{link} and was downloaded on May 13, 2025 for all results presented.} which includes results are from AIONER \cite{luo_aioner_2023} for Named Entity Recognition and GNorm2 \cite{wei_gnorm2_2023} for normalization. In order to construct the MeSH term to Pubmed co-occurrences, we downloaded the full Pubmed Baseline containing all pubmed information for 2025\footnote{The Pubmed baseline XML files are available at the following \href{https://ftp.ncbi.nlm.nih.gov/pubmed/baseline/}{link} and the January 10, 2025 release was used for all results presented.}. From these set of co-occurrences, we then created a master gene set library that would have included publications that were based on genetics including GWAS. Therefore, we curated a list of MeSH terms related to GWAS as seen in Table~\ref{tab:gwas-mesh-terms}. We then formed three MeSH-based gene set libraries: (1) a library in which all pubmed publications are considered, (2) a library in which all GWAS and genetics-based papers are excluded (no circularity), and (3) a library that contains only GWAS and genetics-based papers (fully circular). We then used these three libraries to study the effect of circularity explicitly. Further, we weighted each gene in their associated geneset by a normalized jaccard index and which each MeSH term has an associated gene set. MeSH terms that had more than 2000 genes associated with them were excluded.

\begin{table}[h]
\centering
\begin{tabular}{ll}
\toprule
\textbf{MeSH ID} & \textbf{Concept Name} \\
\midrule
\href{https://meshb.nlm.nih.gov/record/ui?ui=D055106}{D055106} & Genome-Wide Association Study \\
\href{https://meshb.nlm.nih.gov/record/ui?ui=D056726}{D056726} & Genetic Association Studies \\
\href{https://meshb.nlm.nih.gov/record/ui?ui=D000073336}{D000073336} & Whole Genome Sequencing \\
\href{https://meshb.nlm.nih.gov/record/ui?ui=D000073359}{D000073359} & Whole Exome Sequencing \\
\href{https://meshb.nlm.nih.gov/record/ui?ui=D016022}{D016022} & Case-Control Studies \\
\href{https://meshb.nlm.nih.gov/record/ui?ui=D015894}{D015894} & Genome, Human \\
\href{https://meshb.nlm.nih.gov/record/ui?ui=D015810}{D015810} & Linkage Disequilibrium \\
\href{https://meshb.nlm.nih.gov/record/ui?ui=D020022}{D020022} & Genetic Predisposition to Disease \\
\href{https://meshb.nlm.nih.gov/record/ui?ui=D020641}{D020641} & Polymorphism, Single Nucleotide \\
\href{https://meshb.nlm.nih.gov/record/ui?ui=D005838}{D005838} & Genotype \\
\href{https://meshb.nlm.nih.gov/record/ui?ui=D005787}{D005787} & Gene Frequency \\
\href{https://meshb.nlm.nih.gov/record/ui?ui=D040641}{D040641} & Quantitative Trait Loci \\
\href{https://meshb.nlm.nih.gov/record/ui?ui=D023281}{D023281} & Genomics \\
\href{https://meshb.nlm.nih.gov/record/ui?ui=D057182}{D057182} & Mendelian Randomization Analysis \\
\href{https://meshb.nlm.nih.gov/record/ui?ui=D059014}{D059014} & High-Throughput Nucleotide Sequencing \\
\href{https://meshb.nlm.nih.gov/record/ui?ui=D056915}{D056915} & DNA Copy Number Variations \\
\href{https://meshb.nlm.nih.gov/record/ui?ui=D052707}{D052707} & Genomic Instability \\
\href{https://meshb.nlm.nih.gov/record/ui?ui=D004252}{D004252} & DNA Mutational Analysis \\
\href{https://meshb.nlm.nih.gov/record/ui?ui=D008040}{D008040} & Genetic Linkage \\
\href{https://meshb.nlm.nih.gov/record/ui?ui=D019656}{D019656} & Loss of Heterozygosity \\
\bottomrule
\end{tabular}
\caption{MeSH terms excluded to avoid literature derived from GWAS, variant, and sequencing-based genetic analyses.}
\label{tab:gwas-mesh-terms}
\end{table}

\subsection{Relative Success as a Metric}\label{sec:methods-rs}

% ANALYSIS IDEA: Consider adding a supplementary figure showing the sensitivity of
% EA-RS and MLRS rankings to the choice of Laplace smoothing parameter, as this
% could affect the ordering of methods at extreme thresholds.

\paragraph{Discrete relative success.}
To evaluate whether gene-trait prioritization resources enrich for successful drug targets, we adopted the relative success (RS) framework introduced by Minikel et al.\ \cite{minikel_refining_2024}. RS quantifies the degree to which targets designated as ``genetically supported'' advance through clinical trial phases at higher rates than unsupported targets. Drug development programs are tracked across five ordered clinical stages: Preclinical ($j{=}0$), Phase~I ($j{=}1$), Phase~II ($j{=}2$), Phase~III ($j{=}3$), and Launched ($j{=}4$). For each target--indication (T-I) pair, a binary label $G \in \{0,1\}$ indicates whether the pair is classified as genetically supported by a given resource. For each transition $j \to j{+}1$ (where $j = 0, \ldots, 3$), the phase-specific RS is defined as:
\begin{equation}\label{eq:rs-phase}
    \text{RS}_j = \frac{P(\text{advance beyond phase } j \mid G=1)}{P(\text{advance beyond phase } j \mid G=0)} = \frac{X_{G,j} / N_{G,j}}{X_{\bar{G},j} / N_{\bar{G},j}},
\end{equation}
where $N_{G,j}$ and $N_{\bar{G},j}$ are the number of T-I pairs with and without genetic support that reached at least phase $j$, and $X_{G,j}$ and $X_{\bar{G},j}$ are the number that advanced beyond phase $j$ in each group, respectively. The cumulative RS across all phase transitions is then given by the product:
\begin{equation}\label{eq:rs-cum}
    \text{RS}_{\text{cum}} = \prod_{j=0}^{3} \text{RS}_j.
\end{equation}
An $\text{RS}_{\text{cum}} > 1$ indicates that genetically supported targets are more likely to progress through the full clinical pipeline than unsupported targets. Confidence intervals for $\text{RS}_{\text{cum}}$ are computed via the delta method on the log scale: $\text{Var}(\log \text{RS}_j) \approx (1-\hat{p}_{G,j})/(N_{G,j}\hat{p}_{G,j}) + (1-\hat{p}_{\bar{G},j})/(N_{\bar{G},j}\hat{p}_{\bar{G},j})$, where $\hat{p}_{G,j} = X_{G,j}/N_{G,j}$ and $\hat{p}_{\bar{G},j} = X_{\bar{G},j}/N_{\bar{G},j}$. The total variance $\text{Var}(\log \text{RS}_{\text{cum}}) = \sum_j \text{Var}(\log \text{RS}_j)$ yields a $95\%$ confidence interval of $\exp(\log \text{RS}_{\text{cum}} \pm z_{0.975}\sqrt{\text{Var}(\log \text{RS}_{\text{cum}})})$ \cite{minikel_refining_2024}.

\paragraph{Continuous extension for scored resources.}
The standard RS framework requires a binary partition of T-I pairs into supported and unsupported sets, which is natural for resources that provide a discrete classification (e.g., Open Targets Genetics locus-to-gene scores above a fixed threshold) but does not directly accommodate resources that assign continuous scores to gene--trait associations, such as PIGEAN posterior probabilities, MAGMA $p$-values, or Open Targets Platform association scores. To evaluate such resources, we developed a continuous RS curve that traces $\text{RS}_{\text{cum}}$ as a function of the selection fraction $\alpha \in (0, 1]$.

% TODO: CONTINUOUS_RS_N_TI_PAIRS - Total number of T-I pairs used in the continuous RS calculation.
% This would help ground the alpha fractions in absolute numbers.

Concretely, all T-I pairs are sorted by their score in descending order (or ascending order for metrics where lower values indicate stronger evidence, e.g.\ $p$-values). We then sweep a threshold from the most stringent to the most permissive score. At each threshold, the top-$\alpha$ fraction of scored T-I pairs is designated as genetically supported ($G{=}1$) and the remainder---comprising both lower-scoring pairs and pairs absent from the resource entirely---form the unsupported set ($G{=}0$). To maintain numerical stability when subgroups are small, we apply Laplace smoothing to the per-phase success proportions: $\hat{p}_{G,j} = (X_{G,j} + 0.5)/(N_{G,j} + 1)$ and analogously for $\hat{p}_{\bar{G},j}$. This incremental procedure is implemented efficiently using a single pass through the sorted list: as each T-I pair (or block of tied scores) is moved from the unsupported to the supported set, the phase-specific count arrays $N_{G,j}$, $N_{\bar{G},j}$, $X_{G,j}$, and $X_{\bar{G},j}$ are updated in $O(K)$ time per step, where $K{=}4$ is the number of phase transitions, yielding a total time complexity of $O(n \log n + nK)$ for $n$ T-I pairs. At each step, $\text{RS}_{\text{cum}}(\alpha)$ and its delta method confidence interval are recorded.

\paragraph{Summary metrics for continuous RS curves.}
To compare continuous resources on a single scalar, we compute three summary statistics from the RS curve. First, the \textit{excess area above RS\,=\,1} (EA-RS) is defined as the signed area between the RS curve and the baseline $\text{RS}{=}1$:
\begin{equation}\label{eq:ea-rs}
    \text{EA-RS} = \int_0^1 \bigl(\text{RS}_{\text{cum}}(\alpha) - 1\bigr) \, d\alpha,
\end{equation}
computed via trapezoidal integration with an anchor point at $(\alpha{=}0, \text{RS}{=}1)$. EA-RS is positive when the resource provides net enrichment for successful targets across all thresholds and captures both the magnitude and breadth of the enrichment signal. Second, the \textit{mean log RS} (MLRS) is defined as:
\begin{equation}\label{eq:mlrs}
    \text{MLRS} = \int_0^1 \log \text{RS}_{\text{cum}}(\alpha) \, d\alpha,
\end{equation}
which provides a symmetric measure on the log scale in which doubling and halving of RS carry equal weight but opposite sign. The quantity $\exp(\text{MLRS})$ gives the geometric mean RS across the curve. Third, we report $\text{RS}_{\max} = \max_\alpha \text{RS}_{\text{cum}}(\alpha)$, the peak RS achieved at any threshold, which characterizes the best-case enrichment for a given resource. Together, EA-RS captures the overall predictive value integrated across all possible thresholds, MLRS provides a scale-symmetric summary, and $\text{RS}_{\max}$ indicates the ceiling performance of each resource.

\paragraph{Drug target--indication data and matching.}
Clinical trial outcomes were obtained from the Pharmaprojects database, which records drug development programs as target--indication (T-I) pairs annotated with clinical phase information (Preclinical through Launched) and phase transition success indicators \cite{minikel_refining_2024}. Each T-I pair is identified by a unique identifier (\texttt{ti\_uid}) comprising the target gene symbol and the indication MeSH ID. To link gene--trait associations from PIGEAN and other resources to Pharmaprojects T-I pairs, we mapped trait identifiers to MeSH terms and computed pairwise semantic similarities between the MeSH ID assigned to each Pharmaprojects indication and the MeSH ID(s) associated with each gene--trait prediction. When an association mapped to multiple MeSH terms, we retained the maximum similarity. A T-I pair was classified as genetically supported if the best-matching association for that target--indication combination exceeded a similarity threshold of $0.8$, following Minikel et al.\ \cite{minikel_refining_2024}. For the continuous RS extension, this similarity-based matching determines which T-I pairs enter the scored set; pairs whose best similarity falls below the threshold are included in the unsupported background regardless of their score.

% TODO: RS_SIMILARITY_MATRIX_SOURCE - Describe the source and construction of the MeSH similarity matrix.
% This is important for reproducibility -- was it from Minikel et al.? Medical Subject Headings hierarchy?

\paragraph{Target--indication deduplication.}
Because a single T-I pair may appear in multiple rows (e.g., from multiple association sources or MeSH mappings), we deduplicated records before computing RS. For each \texttt{ti\_uid}, we retained the row with the highest MeSH similarity and the most advanced clinical phase. Among remaining ties, we prioritized records with the most complete phase transition outcome data (i.e., the highest phase at which a success or failure indicator was recorded). When computing continuous RS curves for a scored resource, if multiple rows still remained for a given \texttt{ti\_uid} after the above steps, we applied a score aggregation strategy (maximum score by default) to produce a single score per T-I pair.

\paragraph{Benchmarked resources.}
We compared PIGEAN indirect and combined support against three external resources: Open Targets Genetics (OTG) locus-to-gene (L2G) scores, the Open Targets Platform (OTP) harmonic mean genetic association score, and OTP's literature- and model-based indirect association metric. For OTG, we used the L2G score curated by Minikel et al.\ \cite{minikel_refining_2024} with a standard threshold of $0.1$ for the discrete RS comparison. All resources were evaluated using the continuous RS framework described above, with each resource's score column used for ranking (L2G score for OTG; \texttt{overall} and \texttt{drug} scores for OTP; combined, indirect, and direct posterior probabilities for PIGEAN).

% TODO: OTG_L2G_VERSION - Specify the version/date of the OTG L2G data used.
% TODO: OTP_VERSION - Specify the version/date of the OTP data used.

\paragraph{Temporal holdout analysis.}
To assess whether the observed enrichment is driven by circularity---i.e., gene set annotations that directly encode existing drug target knowledge---we performed a temporal holdout analysis. We re-ran PIGEAN using only gene set libraries curated from data available before 2018, thereby excluding any annotations derived from post-2018 drug approvals or clinical trial outcomes. We then restricted the evaluation set to T-I pairs corresponding to drugs launched after 2018, ensuring that the gene set annotations used for scoring could not have been influenced by the clinical outcomes being predicted. RS curves and summary metrics were computed identically to the primary analysis.

\subsection{Common Disease GWAS}
% TODO: Discuss GWAS Catalog curation and portal curation. Discuss mapping to MESH.
Discuss GWAS Catalog curation and portal curation. Discuss mapping to MESH.

\subsection{Rare Disease Curation}
% TODO: Discuss running PIGEAN on gene lists. Discuss mapping to MeSH.
Discuss running PIGEAN on gene lists. Discuss mapping to MeSH.

\subsection{Gene Pleiotropy Quantification}

We quantified gene-level pleiotropy by aggregating PIGEAN Bayes factor scores across 3,026 traits (from 3,077 total; 51 excluded due to insufficient data) for 18,325 genes. For each gene, we computed the weighted sum of positive log Bayes factors across all phenotypes within a trait group. We evaluated three types of PIGEAN evidence: combined evidence, which integrates direct and indirect components; direct evidence (log Bayes factor from the PIGEAN association model); and indirect evidence (prior probability derived from gene set membership). Traits were organized into four groups: all (all 3,026 available traits, serving as the primary analysis), portal (common disease traits from the Type 2 Diabetes Knowledge Portal), GWAS catalog (traits from the NHGRI-EBI GWAS Catalog), and rare (rare disease phenotypes).

\paragraph{Trait correlation correction.}
Because many phenotypes share genetic architecture, summing raw Bayes factors across traits inflates pleiotropy estimates for genes associated with clusters of correlated traits. To address this, we computed pairwise Pearson correlation matrices from posterior-transformed gene-level score profiles (4,576,825 trait pairs) and derived per-trait weights to down-weight redundant phenotypes. Our primary method was inverse correlation weighting, where each trait receives weight $w_t = 1 / \sum_{t'} |C_{t,t'}|$. Weights were normalized to a mean of 1.0 (standard deviation 0.73--0.82 depending on evidence type), with individual weights ranging from 0.16 to 18.0, reflecting substantial variation in trait redundancy.

The effective number of independent traits---estimated as the number of eigenvalues needed to explain the total variance assuming equal contributions---was 130 for combined evidence, 311 for direct evidence, and 61 for indirect evidence (out of 3,026 total traits). The higher effective count for direct evidence (ratio 0.107) compared to indirect evidence (ratio 0.020) indicates that pathway-mediated scores exhibit substantially more inter-trait correlation, consistent with the expectation that many phenotypes share underlying pathway biology. Trait correlations were computed from direct Bayes factors in the primary analysis, as these capture association strength without pathway-derived signal that could introduce additional correlation structure. Alternative correction methods are described in the Supplementary Information.

\subsection{Evolutionary Constraint Metrics}

Gene-level evolutionary constraint metrics were obtained from the gnomAD v4 database. For each gene, we selected the MANE Select transcript (or canonical transcript if unavailable), retaining one transcript per gene. We analyzed six constraint features: LOEUF (loss-of-function observed/expected upper bound confidence interval; lower values indicate greater constraint), lof\_z (loss-of-function z-score), mis\_z (missense z-score), syn\_z (synonymous z-score, serving as a negative control), lof\_oe (loss-of-function observed/expected ratio), and mis\_oe (missense observed/expected ratio). pLI was excluded due to its bimodal distribution (concentrated at 0 and 1), which violates assumptions of linear regression.

\subsection{Pleiotropy--Constraint Regression}

We assessed the relationship between pleiotropy and constraint using univariate linear regression and rank correlation. For each combination of evidence type (combined, direct, indirect), trait group (all, portal, GWAS catalog, rare), and correlation evidence type (combined, direct, indirect), we regressed the weighted pleiotropy sum on each constraint feature individually, yielding 216 regressions across the full parameter grid. We report both Pearson and Spearman rank correlations; Spearman correlations serve as the primary effect size metric due to their robustness to the right-skewed distribution of pleiotropy scores. All regressions included approximately 15,200--16,000 genes depending on constraint feature availability (from 18,325 genes with pleiotropy data, reduced by the inner join with gnomAD constraint annotations).
